\section{Part Three: Define the Job Before Decomposing the Tool}

A product is suitable, or not, only for a specific combination of task, data sensitivity, degree of reliance, and consequence of error. Skipping this leads to arguing about whether a tool is "good" --- a question with no stable answer --- instead of whether it's right for a specific job.

Defining the job means stating, in one sentence: who will use it, for what task, on what data, how the output will be relied upon, and what harm follows if it's wrong, incomplete, or improperly disclosed. The same product can be an easy yes for internal brainstorming on non-sensitive data, a yes-with-mandatory-review for first-pass drafting, and a flat no for autonomous filing --- three verdicts for one product, because they're three different jobs.

Every assessment should terminate in a conditional decision --- approve, approve-with-controls, pilot-only, or reject --- with named owners for residual risk and an explicit list of what remains unknown.
